\section{Contract-Based Assurance in Robotics}
\label{sec:related-work}

The reviewed literature is organized by its \emph{protected object}: the
property or artifact that each approach is designed to preserve.  The classes
overlap, and one system can contribute to more than one class.

\subsection{Related Surveys and Scope}

Recent surveys use several organizing lenses.  Street \emph{et al.} review formal models for robot teams
under uncertainty, while Gielis \emph{et al.} analyze the interaction between
robotic applications and communication systems
\cite{street2023formal,gielis2022communications}.  Garg \emph{et al.} organize
learning-based safe control by safety notions, control methods, and verification
tools \cite{garg2024safecontrol}.

Those surveys clarify capabilities, models, communication, or intervention
layers.  The present survey uses a different unit of comparison: the protected
object and the information required when its guarantee becomes another
component's assumption.  This narrower lens motivates the interface-focused
taxonomy and reference framework below.

\subsection{Mission and Intent Contracts}

CROME represents contextual robotic missions as hierarchical assume--guarantee
contracts and synthesizes sub-mission controllers
\cite{mallozzi2020crome,mallozzi2023contextual}.  Contract algebra can also
refine a mission against a component library or compute a repair when the
original specification is unrealizable \cite{mallozzi2022repair}.  These methods
make goals, contexts, refinement, and realizability explicit at the mission
boundary.  They focus on the mission specification; data age, provenance, and
implementation identity remain outside their scope.

Language-conditioned multi-robot planners address an adjacent problem: turning
open-ended instructions into plans, role assignments, or optimization models.
RoCo uses dialogue among language-model agents, while LiP-LLM combines a
language model with a predefined skill set, dependency graph, and linear
program \cite{mandi2023roco,obata2024lipllm}.  These systems provide planning
capability evidence, but their collision models, skill definitions, and
feasibility solvers remain external assurance mechanisms.  Expressing an
intent is therefore not equivalent to certifying its physical or distributed
execution.

\subsection{Quantitative Assurance of Coordinated Autonomy}

Chaki, Dolan, and Giampapa frame assurance of coordinated autonomy as the
management of quantitative performance expectations.  Their approach combines
a characterization of coordination phenomena, phases of a coordinated mission,
and discrete-time Markov-chain models evaluated through probabilistic model
checking \cite{chaki2013assurance}.  Related work formalizes multi-robot mission
utility as temporal or reward properties and derives compositional results for
probabilistic verification \cite{chaki2013probabilistic}.

Complementary reliability studies estimate the performance of sequentially
coordinated missions from measured atomistic robot behavior, or compare
repairable robot-team configurations at the mission level
\cite{porter2013reliability,stancliff2006missionreliability}.  Within the
reviewed corpus, this line is the closest mission-level predecessor to the
present survey's multi-robot assurance question: it connects coordination
structure to quantitative mission evidence.  Its primary object is mission or
team reliability, whereas the present coverage analysis asks which semantic,
physical, temporal, compatibility, provenance, lifecycle, and recovery
obligations must cross the heterogeneous interfaces that realize such a
mission.  The two views are complementary.

\subsection{Software-Component Contracts}

Luckcuck \emph{et al.} specify ROS nodes with first-order-logic contracts,
derive system properties with composition rules, and translate a ROS Contract
Language into executable monitors \cite{luckcuck2022compositional}.  This work
connects behavioral contracts to node implementations and demonstrates that
component assumptions can be made executable.  Its primary object is node
behavior; cross-robot state age, team membership, and heterogeneous deployment
compatibility are not the central object of the calculus.

ROSMonitoring~2.0 instruments topics and services and monitors event order,
including partial ROS~2 support \cite{ghaffari2024rosmonitoring}.  Caldas
\emph{et al.} complement specification work with guidelines for runtime
verification and field-based testing, emphasizing observability,
instrumentation, environment, and reporting \cite{caldas2024rvft}.  These
approaches show that specification, instrumentation, and evidence collection
are separate engineering concerns.

\subsection{Distributed Multi-Robot Programming and Design}

Koord provides platform-independent abstractions for distributed shared state,
robot coordination, sensing, and control.  Its executable semantics expose
proof obligations that can be discharged using symbolic and physics-based
verification tools \cite{ghosh2020koord}.  ATLAS takes a model-driven approach:
it represents team architecture, mission requirements, and collective
intelligence, then uses simulation to explore design choices and robustness
\cite{harbin2023design}.  Both connect software structure to multi-robot
behavior more directly than a message-level monitor.  Their analyses remain
tied to their languages, models, generated infrastructure, and declared
simulation settings; they do not define a general evidence envelope for every
interface in a heterogeneous deployed stack.

\subsection{Dynamical and Multi-Agent Contracts}

Contracts for dynamical systems place assumptions on input trajectories and
guarantees on outputs.  Sharf \emph{et al.} provide satisfaction and refinement
tools for linear contracts and illustrate them with two vehicles
\cite{sharf2021dynamical}.  Chen \emph{et al.} use parameterized temporal-logic
contracts to compute invariant sets compositionally in networked systems
\cite{chen2019invariance}.  Liu \emph{et al.} synthesize decentralized
controllers for interconnected nonlinear systems from signal-temporal-logic
contracts \cite{liu2022stlcontracts}.

Assume--guarantee reinforcement learning instead represents neighboring-agent
behavior through regular-language abstractions and combines local satisfaction
probabilities into a global lower bound \cite{kazemi2024agrl}.  These approaches
directly address interaction and decomposition, but typically assume that the
declared neighbor abstraction, labels, communication observations, dynamics,
and timing semantics are available as modeled.

\subsection{Distributed State and Communication}

ROS~2 provides modular robotics communication commonly through DDS
\cite{macenski2022ros2,omg2015dds}.  DDS quality-of-service policies govern such
properties as reliability, history, durability, liveliness, deadline, and
lifespan.  These transport and data-stream mechanisms do not prove
application-level freshness or mission validity.  For example, the DDS
\emph{Deadline} policy concerns expected sample cadence; it does not establish
the age of the physical event represented by an arriving sample.

ROS~2 latency studies show that process structure, executor behavior, and DDS
implementation affect end-to-end timing
\cite{maruyama2016ros2,kronauer2021latency}.  Communication surveys likewise
show that coordination and networking assumptions are often designed and
evaluated separately \cite{gielis2022communications}.  Recent planning work
incorporates temporal communication constraints and intermittent connectivity
directly into stochastic multi-robot planning \cite{silva2026probabilistic}.
Age-of-information research
separately characterizes how fresh the newest available update is
\cite{kaul2012aoi}.  These works provide deployment semantics and measurements,
not a compositional robot contract.  They nevertheless expose a necessary
bridge: a dynamical or safety guarantee that assumes current peer state needs an
application-level rule relating event time, clock uncertainty, delivery, and
the physical decision.

\subsection{Runtime Assurance and Safety}

SOTER pairs an advanced controller with a safety controller and a decision
module that transfers authority when a monitor predicts a violation
\cite{desai2019soter}.  SOTER on ROS evaluates this pattern in simulated
multi-robot surveillance \cite{shivakumar2020soterros}.  Control barrier
functions and multi-robot barrier certificates instead constrain the admissible
control set to preserve a modeled safe set
\cite{ames2017cbf,wang2017barrier}.  Both approaches provide strong conditional
safety mechanisms.  Their guarantees still depend on dynamics, state
estimation, update rate, feasibility, and the prompt availability of the
required peer state.  A local intervention can also preserve one robot's safety
while invalidating a team plan unless the intervention is reported upstream.

\subsection{Cross-Layer Guardrails}

Modular safety guardrails motivate monitoring and intervention across action,
decision, and human-centered safety \cite{kim2026guardrails}.  The position
leaves open how monitors in different layers and robots share evidence.
Within the reviewed works, I did not find one composition model that jointly
represents temporal validity, ownership, provenance, and recovery across
cross-layer and peer-robot interfaces.

\begin{table}[t]
  \centering
  \caption{Cross-layer taxonomy of selected assurance approaches}
  \label{tab:related-work}
  \scriptsize
  \setlength{\tabcolsep}{1.4pt}
  \renewcommand{\arraystretch}{1.03}
  \begin{tabularx}{\textwidth}{@{}P{0.105\textwidth}P{0.16\textwidth}P{0.14\textwidth}P{0.21\textwidth}P{0.16\textwidth}Y@{}}
    \toprule
    Class & Representative work & Protected object & Mechanism and assurance evidence & Deployment / MRS evidence & Underemphasized or external obligations \\
    \midrule
    Mission and intent & CROME; contextual missions; specification repair \cite{mallozzi2020crome,mallozzi2023contextual,mallozzi2022repair} & Goals, contexts, and component libraries & Contract algebra, synthesis, refinement, switching, and repair & Mission examples and synthesized controllers; no physical team evaluation & Physical validity, event-time state, implementation identity, and peer reconciliation \\
    Coordinated-autonomy assurance & Quantitative assurance, probabilistic mission verification, and reliability estimation \cite{chaki2013assurance,chaki2013probabilistic,porter2013reliability,stancliff2006missionreliability} & Mission phases, coordination commitments, utility, and team reliability & DTMC model checking, compositional probabilistic verification, and conditional performance estimation & Initial robot-validation results, atomistic behavior measurements, and repairable-team analyses & Interface-carried uncertainty, provenance, software identity, lifecycle state, and recovery evidence \\
    Learned task planning & RoCo; LiP-LLM \cite{mandi2023roco,obata2024lipllm} & Language-conditioned tasks, roles, dependencies, and plans & Environment feedback or optimization provides external feasibility checks & Multi-robot benchmarks, simulation, and selected real-world execution & Calibrated uncertainty, timing, compatibility, and deployment assurance \\
    Component behavior & ROS contracts \cite{luckcuck2026compositional} & ROS node assumptions and guarantees & First-order contracts, inference rules, RCL, and generated runtime monitors & Rover and component examples; no robot-team study & Team epoch, ownership, state age, and evidence lineage \\
    Distributed programming and design & Koord; ATLAS \cite{ghosh2020koord,harbin2023design} & Coordination code, proof obligations, team architecture, mission, and collective intelligence & Executable semantics, modular verification, generated infrastructure, simulation, and robustness analysis & Three Koord applications and three simulated ATLAS case studies & Runtime freshness, evidence lineage, heterogeneous implementation identity, and recovery authority \\
    Dynamical systems & Linear and temporal contracts \cite{sharf2021dynamical,chen2019invariance,liu2022stlcontracts} & Input/output trajectories, temporal tasks, and coupled dynamics & Satisfaction, refinement, invariant-set computation, and decentralized synthesis & Two-vehicle and interconnected-system simulations & Runtime message age, lifecycle, implementation, and recovery semantics \\
    Multi-agent learning & Assume--guarantee RL \cite{kazemi2024agrl} & Neighbor abstractions, labels, and local policies & Local learning with a lower bound on global satisfaction probability & Grid-world, building, and traffic-control studies & Physical and middleware validation of the assumed abstractions \\
    Distributed state & ROS~2, DDS, communication studies, and age of information \cite{macenski2022ros2,omg2015dds,gielis2022communications,kronauer2021latency,kaul2012aoi,silva2026probabilistic} & Data-stream semantics, delivery status, freshness, and intermittent connectivity & QoS compatibility/status, timing measurements, and communication-aware planning & Standards, multi-node benchmarks, and multi-robot planning simulations & Physical decision validity, evidence lineage, and coordinated recovery \\
    Runtime assurance & ROSMonitoring; SOTER; barrier certificates; field-testing and safe-control surveys \cite{ghaffari2024rosmonitoring,shivakumar2020soterros,wang2017barrier,caldas2024rvft,garg2024safecontrol} & Events, safety properties, and controller authority & Event monitors, safe-control sets, controller switching, instrumentation, and reporting & UAV case, simulated surveillance, physical multi-robot safety, and practitioner evidence & Semantic evidence, implementation compatibility, and upstream intervention feedback \\
    Cross-layer guardrails and runtime & Guardrails; PhyAgentOS \cite{kim2026guardrails,liu2026phyagentos} & Layered interventions, execution sessions, and evidence & Monitoring, preflight, verification, safety, heartbeat, and evidence-grounded completion & Simulated and physical execution across embodiments & Common composition semantics across heterogeneous peer-robot edges \\
    \bottomrule
  \end{tabularx}
\end{table}


\subsection{Survey Synthesis}

The reviewed literature provides mechanisms for mission realizability, node
behavior, dynamical composition, communication observation, and runtime safety.
Coverage becomes discontinuous where these mechanisms meet.  A mission
guarantee may end before physical
feasibility; a dynamics guarantee may assume current neighbor signals; a
middleware status may omit the physical event time; and a safety intervention
may not return to task allocation.  \Cref{sec:gaps} converts these discontinuities
into a compact set of cross-layer coverage gaps.
